%=============================================================================
% WAVE 4, ITERATION 10: EXPERIMENTAL PROPOSALS
% Four Actionable Mini-Proposals for DFD Experimental Tests
%
% Agent 10 | DFD Research Programme | Model: Opus 4.6
% Building on: MASTER_FINDINGS_CORPUS.md, W3i7_P1_P2_P11_update.tex,
%              W3i8_Detection_update.tex, section02_formalism.tex
%
% Status flags: [VERIFIED], [CONJECTURE], [NEW], [ACTIONABLE]
% Date: 20 March 2026
%=============================================================================

\documentclass[11pt]{article}
\usepackage[utf8]{inputenc}
\usepackage{amsmath,amssymb,amsfonts,mathtools}
\usepackage{geometry}
\geometry{margin=1in}
\usepackage{booktabs}
\usepackage{siunitx}
\usepackage{hyperref}
\usepackage{xcolor}
\usepackage{enumitem}
\usepackage{float}
\usepackage{array}
\usepackage{longtable}
\usepackage{tcolorbox}
\tcbuselibrary{skins,breakable}

\newcommand{\ee}[1]{\times 10^{#1}}
\newcommand{\lbone}{|\lambda - 1|}
\newcommand{\dpsi}{\delta\psi}
\newcommand{\Qpsi}{Q_\psi}
\newcommand{\SNR}{\mathrm{SNR}}

\definecolor{proposalcolor}{rgb}{0.0,0.35,0.55}
\definecolor{actioncolor}{rgb}{0.0,0.5,0.0}

\title{Wave 4, Iteration 10: Four Actionable Experimental Proposals\\
\large LIGO Archival Reanalysis $\cdot$ GPS 59\,ps Archival Search\\
Si$_3$N$_4$ MEMS Prototype $\cdot$ Galaxy Rotation Curve Reanalysis\\[6pt]
\normalsize Density Field Dynamics Research Programme}

\author{DFD Research Programme --- Experimental Proposals Agent (W4i10)\\
\textit{Building on: W3i1--W3i9 corpus, MASTER\_FINDINGS\_CORPUS}}

\date{20 March 2026}

\begin{document}
\maketitle

%=============================================================================
\section*{Executive Summary}
%=============================================================================

\begin{tcolorbox}[colback=blue!5!white, colframe=proposalcolor,
  title=\textbf{Four Proposals: From Zero-Cost to \$500K}]

This document contains four self-contained experimental proposals,
each written so that a scientist can begin work immediately.
All technical parameters trace back to the converged DFD corpus
(W3i8 final, 60+ corrections applied).

\begin{center}
\renewcommand{\arraystretch}{1.3}
\begin{tabular}{@{}clccc@{}}
\toprule
\textbf{\#} & \textbf{Proposal} & \textbf{Budget} & \textbf{Timeline} & \textbf{DFD Reach} \\
\midrule
1 & LIGO O3 Archival Reanalysis & \$100--500K & 12 months & $\lbone \sim 3\ee{-14}$ \\
2 & GPS 59\,ps Archival Search & \$50--100K & 6 months & SNR $>100$ \\
3 & Si$_3$N$_4$ MEMS Prototype & \$200--500K & 12--18 months & $Q \geq 10^9$ at 20\,mK \\
4 & Galaxy Rotation Curve Reanalysis & \$0 & 3--6 months & $>5\sigma$ model selection \\
\bottomrule
\end{tabular}
\end{center}

Each proposal specifies: data sources, analysis methods, systematics
budget, null tests, personnel, timeline, and success criteria.
\end{tcolorbox}

\tableofcontents
\newpage


%=============================================================================
%=============================================================================
\part{Proposal 1: LIGO O3 Archival Reanalysis}
\label{part:ligo}
%=============================================================================
%=============================================================================

\begin{tcolorbox}[colback=green!3!white, colframe=actioncolor,
  title=\textbf{PROPOSAL 1 --- LIGO ARCHIVAL REANALYSIS}]
\textbf{Budget}: \$100--500K (postdoc salaries + computing)\\
\textbf{Timeline}: 12 months\\
\textbf{DFD reach}: $\lbone \sim 3\ee{-14}$ (gradient channel at 100\,Hz)\\
\textbf{Data}: Publicly available O1/O2/O3 strain data from GWOSC\\
\textbf{Personnel}: 1 PI (20\% FTE) + 1--2 postdocs (100\% FTE)
\end{tcolorbox}

%=============================================================================
\section{Scientific Objective}
%=============================================================================

Search for a scalar breathing mode $h_\psi = 2\,\dpsi$ in LIGO/Virgo
archival data. DFD predicts that EM energy density sources a scalar
perturbation $\dpsi$ that propagates at $c$ (massless, non-dispersive)
and couples to LIGO through the gradient channel with suppression
$\sin^2(\pi L/\lambda_\psi)$ relative to tensor modes (W3i7, Theorem 3.1).

A detection or stringent upper limit at $\lbone \sim 3\ee{-14}$
would improve the current authoritative SQMS bound
($\lbone < 1.56\ee{-9}$, W3i3) by five orders of magnitude.

%=============================================================================
\section{Data Specification}
%=============================================================================

\subsection{Source: Gravitational Wave Open Science Center (GWOSC)}

All strain data are publicly available at \url{https://gwosc.org/data/}.

\begin{center}
\renewcommand{\arraystretch}{1.25}
\begin{tabular}{@{}llllc@{}}
\toprule
\textbf{Run} & \textbf{Dates} & \textbf{Detectors} & \textbf{Sample Rate} & \textbf{Livetime} \\
\midrule
O1 & Sep 2015 -- Jan 2016 & H1, L1 & 4096/16384\,Hz & $\sim$130 days \\
O2 & Nov 2016 -- Aug 2017 & H1, L1, V1 & 4096/16384\,Hz & $\sim$270 days \\
O3a & Apr 2019 -- Oct 2019 & H1, L1, V1 & 4096/16384\,Hz & $\sim$180 days \\
O3b & Nov 2019 -- Mar 2020 & H1, L1, V1 & 4096/16384\,Hz & $\sim$147 days \\
\bottomrule
\end{tabular}
\end{center}

\textbf{Strain channels to use}: \texttt{H1:GDS-CALIB\_STRAIN} and
\texttt{L1:GDS-CALIB\_STRAIN} (calibrated strain, $h(t)$).
For Virgo: \texttt{V1:Hrec\_hoft\_16384Hz}.

\textbf{Data format}: Available as HDF5 and GWF (frame) files.
Use 4096\,Hz data for the primary analysis (adequate for $f < 2000$\,Hz);
16384\,Hz data for high-frequency checks.

\textbf{Auxiliary channels}: Download O3 auxiliary channel data release
(\url{https://gwosc.org/O3/auxiliary/}) for Schumann resonance
magnetometer data (\texttt{H1:PEM-CS\_MAG\_LVEA\_VERTEX\_Z}) and
seismometer channels for Newtonian noise regression.

\subsection{Data quality}

Use the GWOSC data quality flags:
\begin{itemize}[nosep]
\item \texttt{DATA\_QUALITY\_VECTOR}: keep only \texttt{SCIENCE\_MODE}
  segments (bit 0 = 1).
\item \texttt{CBC\_CAT1}: remove hardware injections and known glitch
  periods.
\item Apply standard gating for loud glitches ($> 50\sigma$ in whitened
  strain).
\end{itemize}

Total usable livetime after quality cuts: $\sim$500 days across
O1/O2/O3 combined.

%=============================================================================
\section{Analysis Method}
%=============================================================================

\subsection{Strategy overview}

The scalar breathing mode couples to LIGO through the \textbf{gradient
channel}: spatial variation of $\dpsi$ across the detector arms produces
a differential phase (W3i7 Eq.~3.4):
\begin{equation}
h_{\rm diff}^{\rm scalar}(f) = 2\,\dpsi(f)\,\sin^2\!\left(
\frac{\pi L}{c/f}\right)\cos\theta,
\label{eq:gradient_coupling}
\end{equation}
where $L = 4\,$km (LIGO arm length), $f$ is the signal frequency,
and $\theta$ is the angle between the scalar wavevector and the arm
bisector.

At $f = 100$\,Hz: $\sin^2(\pi L f/c) \sim 1.7\ee{-5}$.
At $f = 500$\,Hz: $\sin^2(\pi L f/c) \sim 4.4\ee{-4}$.

\subsection{Three search channels}

\subsubsection{Channel A: Targeted magnetar flare correlation}

\textbf{Rationale}: DFD predicts $\dpsi$ sourced by EM energy release.
Magnetar giant flares release $10^{44}$--$10^{47}$\,erg in EM energy
over $\sim 0.1$--$1$\,s. This is the strongest impulsive EM event
accessible to LIGO's frequency band.

\textbf{Method}:
\begin{enumerate}[nosep]
\item Compile catalog of magnetar flares during O1/O2/O3 from
  \textit{Fermi}/GBM and \textit{Swift}/BAT trigger lists.
\item For each flare, extract $\pm 10$\,s of strain data from all
  active LIGO/Virgo detectors.
\item Compute the excess cross-power in the scalar polarisation
  channel using the Bayesian polarisation decomposition framework
  (as used by the LVK for GW170817 polarisation tests).
\item The scalar antenna pattern is $F_\psi(\hat{n}) = \frac{1}{2}\sin^2\zeta$,
  where $\zeta$ is the angle between the source direction and the
  detector normal. This differs from the tensor patterns and allows
  polarisation separation with $\geq 3$ detectors.
\end{enumerate}

\textbf{Expected reach}: $\lbone \sim 3\ee{-14}$ for a giant flare
at 10\,kpc (W3i7 Table~3.2).

\subsubsection{Channel B: Continuous-wave solar search}

\textbf{Rationale}: The Sun is the strongest continuous EM source.
DFD predicts a psi-perturbation with diurnal modulation as the
detector rotates relative to the Sun.

\textbf{Method}:
\begin{enumerate}[nosep]
\item Compute the time-dependent scalar antenna pattern for the Sun's
  position as seen from H1 and L1.
\item Apply a matched filter using the predicted diurnal + annual
  amplitude modulation template.
\item Cross-correlate H1 and L1 in the scalar channel, weighting
  by the scalar antenna pattern.
\item Integration over full O3 ($\sim$330 days) provides bandwidth
  narrowing of $\sqrt{T \cdot \Delta f}$.
\end{enumerate}

\textbf{Expected reach}: $\lbone \sim 5\ee{-21}$ (optimistic,
contingent on Schumann subtraction to $<10^{-3}$ level).

\subsubsection{Channel C: Stochastic background cross-correlation}

\textbf{Rationale}: A cosmological psi-background with specific
angular spectrum tracing the EM luminosity distribution.

\textbf{Method}: Standard LIGO stochastic search pipeline
(\texttt{PyStoch}) modified to include scalar overlap reduction
function $\gamma_\psi(f)$ instead of tensor $\gamma_T(f)$.

\textbf{Expected reach}: $\Omega_\psi(f) < 10^{-8}$ at 25\,Hz
with O3 data.

%=============================================================================
\section{Systematic Error Budget}
%=============================================================================

\begin{center}
\renewcommand{\arraystretch}{1.25}
\begin{tabular}{@{}p{4cm}p{3cm}p{2.5cm}p{4cm}@{}}
\toprule
\textbf{Systematic} & \textbf{Magnitude} & \textbf{$f$-dependence} &
  \textbf{Mitigation} \\
\midrule
Newtonian noise & $h \sim 10^{-20}$ at 10\,Hz & $\propto f^{-4}$ &
  Restrict analysis to $f > 30$\,Hz; Wiener filter with seismometer
  array data \\
\midrule
Schumann resonances & $h \sim 10^{-24}$ at 14.1\,Hz harmonics &
  Narrow lines & Magnetometer subtraction; excise $\pm 0.5$\,Hz bands
  around 7.83, 14.1, 20.3, 26.4, 32.6\,Hz \\
\midrule
Laser frequency noise & $< 10^{-7}$\,Hz/$\sqrt{\rm Hz}$ &
  $\propto 1/f$ & Common-mode; mimics scalar. Subtract using
  pre-mode-cleaner error signal \\
\midrule
Calibration uncertainty & 5\% amplitude, $3^\circ$ phase &
  $f$-dependent & Use O3 time-dependent calibration model;
  propagate into posterior \\
\midrule
Correlated magnetic noise & Variable & Broadband &
  Cross-correlate with PEM magnetometers; Wiener subtract \\
\bottomrule
\end{tabular}
\end{center}

\textbf{Dominant systematic}: Newtonian noise below 30\,Hz is
indistinguishable from a scalar psi-perturbation in a single
detector (W3i7 Theorem~3.2). Analysis must be restricted to
$f > 30$\,Hz or use multi-detector polarisation decomposition.

%=============================================================================
\section{Five Null Tests}
\label{sec:ligo_null}
%=============================================================================

These null tests are drawn from W3i7, Section~4:

\begin{enumerate}
\item \textbf{Polarisation separation}: With 3+ non-co-aligned detectors
  (H1, L1, V1 in O3), decompose each candidate into tensor ($+, \times$)
  and scalar ($\psi$) components. A DFD signal is scalar-only with
  isotropic antenna pattern $F_\psi = \frac{1}{2}\sin^2\zeta$; a GR
  signal has the characteristic $\cos 2\alpha$ tensor pattern. Requires
  per-detector $\SNR \gtrsim 8$.

\item \textbf{EM--GW time-domain correlation}: Cross-correlate scalar
  channel excess power with EM event catalogs (magnetar flares, solar
  flares, GRBs). DFD predicts correlation; GR predicts none. For BNS
  mergers, the psi-signal should track post-merger EM emission, not the
  inspiral waveform.

\item \textbf{Dispersion test}: DFD predicts zero dispersion (massless
  scalar, $c_s = c$). Measure group velocity across the LIGO band
  (20--2000\,Hz). Any frequency-dependent arrival time rules out DFD
  and points to a massive scalar alternative.

\item \textbf{Stochastic anisotropy}: The angular power spectrum of a
  psi-background traces EM luminosity weighting, not mass weighting.
  Compare measured $C_\ell^{\rm scalar}$ with predictions from EM
  luminosity catalogs and mass distribution catalogs.

\item \textbf{Day/night modulation}: Solar psi-signal produces a
  diurnal modulation with amplitude following the scalar antenna
  pattern $F_\psi(\hat{n}_\odot(t))$. Compare measured diurnal
  amplitude with the predicted template. Cross-check: the annual
  amplitude should scale as $1/r_\odot^2(t)$ (7\% annual variation),
  while thermal systematics have different seasonal dependence.
\end{enumerate}

%=============================================================================
\section{Personnel, Timeline, and Budget}
%=============================================================================

\begin{center}
\renewcommand{\arraystretch}{1.25}
\begin{tabular}{@{}p{3cm}p{3cm}p{3cm}p{4cm}@{}}
\toprule
\textbf{Role} & \textbf{FTE} & \textbf{Duration} & \textbf{Tasks} \\
\midrule
PI (GW data analysis expert) & 0.2 FTE & 12 months &
  Supervision, systematic budget validation, paper writing \\
Postdoc 1 (data analyst) & 1.0 FTE & 12 months &
  Pipeline development, Channel A + C searches, null tests \\
Postdoc 2 (optional) & 1.0 FTE & 12 months &
  Channel B (solar search), Schumann subtraction, EM catalog
  cross-correlation \\
\bottomrule
\end{tabular}
\end{center}

\textbf{Timeline}:
\begin{itemize}[nosep]
\item Months 1--3: Pipeline development. Implement scalar antenna
  pattern in existing stochastic/CW search codes. Download and
  validate O1/O2/O3 data. Build EM event catalog.
\item Months 4--6: Channel A search (magnetar flare correlation).
  This is the fastest path to a result.
\item Months 7--9: Channel B (solar CW) and Channel C (stochastic).
  Schumann subtraction development.
\item Months 10--12: Null tests, systematic studies, paper preparation.
\end{itemize}

\textbf{Budget breakdown}:
\begin{itemize}[nosep]
\item 1 postdoc salary (12 months): \$80--120K (institution-dependent)
\item Optional 2nd postdoc: \$80--120K
\item PI buyout (0.2 FTE): \$30--60K
\item Computing (LIGO Data Grid or cloud): \$10--30K
\item Travel (LVK collaboration meetings): \$5--10K
\item \textbf{Total}: \$125--340K (1 postdoc) or \$205--500K (2 postdocs)
\end{itemize}

\textbf{Funding targets}: NSF PHY (Gravitational Physics program),
NASA Astrophysics Data Analysis Program (ADAP), DOE HEP.

%=============================================================================
\section{Success Criteria}
%=============================================================================

\begin{itemize}
\item \textbf{Detection}: Scalar excess at $>5\sigma$ in $\geq 2$
  null tests simultaneously. Publication as a discovery paper.
\item \textbf{Upper limit}: $\lbone < 3\ee{-14}$ (improving SQMS
  bound by $5\times$). Publication as a constraint paper.
\item \textbf{Null result}: No excess above systematic floor.
  Still publishable: first scalar mode search in LIGO data with
  DFD-specific templates. Sets the stage for O4/O5 searches.
\end{itemize}


\newpage
%=============================================================================
%=============================================================================
\part{Proposal 2: GPS 59\,ps Archival Search}
\label{part:gps}
%=============================================================================
%=============================================================================

\begin{tcolorbox}[colback=green!3!white, colframe=actioncolor,
  title=\textbf{PROPOSAL 2 --- GPS 59\,PS ARCHIVAL SEARCH}]
\textbf{Budget}: \$50--100K (1 geodesist, 6 months)\\
\textbf{Timeline}: 6 months\\
\textbf{Expected SNR}: $>100$ with 2 years of data\\
\textbf{Data}: IGS precise point positioning products (publicly available)\\
\textbf{Personnel}: 1 geodesist/timing specialist
\end{tcolorbox}

%=============================================================================
\section{Scientific Objective}
%=============================================================================

DFD predicts a 59\,ps peak-to-peak annual variation in GPS satellite
clock residuals arising from the scalar psi-field gradient of the
Sun (W3i7; one of only 2 genuinely archival-testable DFD predictions).
This signal has a characteristic diurnal + annual modulation pattern
that is distinct from known GPS systematics. Detection or exclusion
at $\SNR > 100$ is achievable with existing archival data.

The predicted nonreciprocal timing delay for a GPS satellite at zenith
is $\delta t_{\rm NR} = 0.142$\,ns (W3i8 Table). The annual modulation
due to the Earth's orbital eccentricity and changing Sun--Earth geometry
produces a peak-to-peak variation of 59\,ps, which is the target signal.

%=============================================================================
\section{Data Specification}
%=============================================================================

\subsection{Primary data: IGS Final Clock Products}

The International GNSS Service (IGS) provides precise clock solutions
freely at \url{https://igs.org/products-access/} and via CDDIS
(\url{https://cddis.nasa.gov/}).

\textbf{Products to download}:
\begin{itemize}[nosep]
\item \textbf{IGS Final clock RINEX} (\texttt{*.clk}): 30-second
  cadence satellite and station clock solutions. Latency: 12--18 days.
  Available from 1994--present.
\item \textbf{IGS Final orbits} (\texttt{*.sp3}): Precise satellite
  ephemerides at 15-minute intervals.
\item \textbf{IGS Final troposphere ZPD} (\texttt{*.zpd}): Zenith
  path delay estimates at 5-minute intervals for tropospheric
  correction.
\item \textbf{Station coordinates} (\texttt{*.snx}): SINEX files
  with station positions, velocities, and discontinuities.
\end{itemize}

\textbf{Data volume}: $\sim$2 years of Final products $\approx$ 50\,GB
(compressed). Recommend 2019--2021 (overlapping with O3 for possible
joint analysis).

\subsection{Secondary data: Multi-GNSS products}

For the Multi-GNSS ratio test (smoking gun):
\begin{itemize}[nosep]
\item \textbf{IGS MGEX} (Multi-GNSS Experiment): Clock products for
  GPS, GLONASS, Galileo, and BeiDou. Available at
  \url{https://igs.org/mgex/products/}.
\item The predicted Galileo/GPS timing ratio is
  $\delta t_{\rm Gal}/\delta t_{\rm GPS} = 1.070$ (fixed, independent
  of atmospheric systematics; W3i8 confirmed).
\end{itemize}

%=============================================================================
\section{Analysis Method}
%=============================================================================

\subsection{Step 1: Precise Point Positioning residuals}

\begin{enumerate}[nosep]
\item Process raw RINEX observation files from $\geq 50$ globally
  distributed IGS stations using PPP software (recommended:
  \texttt{GIPSY-OASIS} from JPL or \texttt{Bernese GNSS Software}).
\item Obtain satellite clock residuals $\Delta t_{\rm sat}(t)$ after
  removing: satellite orbit errors, receiver clock offsets,
  ionospheric delay (dual-frequency L1/L2 combination), tropospheric
  delay (VMF3 mapping function + estimated ZPD), solid Earth tides
  (IERS 2010 conventions), ocean tide loading (FES2014b), antenna
  phase center variations (IGS ANTEX files).
\end{enumerate}

\subsection{Step 2: Matched filter for DFD signature}

The DFD signal template is:
\begin{equation}
\delta t_{\rm DFD}(t) = A_0 \left[
  1 + \epsilon\,\cos\!\left(\frac{2\pi t}{T_{\rm year}} + \phi_{\rm annual}\right)
\right]
\cos\!\left(\frac{2\pi t}{T_{\rm sid}} + \phi_{\rm diurnal}(t)\right),
\label{eq:gps_template}
\end{equation}
where:
\begin{itemize}[nosep]
\item $A_0 = 59$\,ps (peak annual amplitude)
\item $\epsilon = 0.034$ (orbital eccentricity modulation)
\item $T_{\rm year} = 365.25$\,days
\item $T_{\rm sid} = 23$\,h\,56\,min (sidereal day)
\item $\phi_{\rm diurnal}(t)$ encodes the satellite geometry and
  Sun position; computed from ephemerides.
\end{itemize}

Apply the matched filter:
\begin{equation}
\SNR = \frac{\sum_i w_i\,\Delta t_i\,s_i}
{\sqrt{\sum_i w_i^2\,\sigma_i^2\,s_i^2}},
\end{equation}
where $s_i$ is the template evaluated at time $t_i$, $\sigma_i$ is
the clock residual uncertainty, and $w_i$ are inverse-variance weights.

\subsection{Step 3: Multi-GNSS ratio test}

\begin{enumerate}[nosep]
\item Repeat Step 1--2 independently for Galileo satellites (altitude
  23,222\,km vs GPS 20,200\,km).
\item Compute the ratio $R = \hat{A}_{\rm Gal}/\hat{A}_{\rm GPS}$.
\item DFD prediction: $R = 1.070$ (fixed).
\item Null (no DFD): $R$ is consistent with unity or with a ratio
  determined by atmospheric/orbital systematics (which have different
  altitude dependence).
\end{enumerate}

This ratio test is the ``smoking gun'' because atmospheric systematics
(troposphere, ionosphere) cancel in the ratio, while the DFD signal
does not.

%=============================================================================
\section{Systematics}
%=============================================================================

\begin{center}
\renewcommand{\arraystretch}{1.25}
\begin{tabular}{@{}p{4cm}p{2.5cm}p{6.5cm}@{}}
\toprule
\textbf{Systematic} & \textbf{Magnitude} & \textbf{Mitigation} \\
\midrule
Tropospheric delay residual & $\sim$1--5\,mm ($\sim$3--17\,ps) &
  VMF3 mapping function + ZPD estimation; ratio test cancels
  troposphere \\
\midrule
Ionospheric higher-order terms & $\sim$0.1--1\,ps &
  Third-order ionospheric correction using L1/L2/L5 triple-frequency
  on modernized GPS and Galileo \\
\midrule
Antenna phase center variations & $\sim$2--5\,mm &
  IGS ANTEX calibration tables; satellite-specific corrections \\
\midrule
Relativistic modeling errors & $\sim$0.01\,ps &
  Full IERS 2010 relativistic model; negligible \\
\midrule
Satellite attitude/yaw steering & $\sim$0.1--1\,ps &
  Use only nominal-yaw periods; flag eclipse seasons \\
\midrule
Reference frame drift & $\sim$0.1\,mm/yr &
  Use IGS Final frame realization; negligible over 2\,yr \\
\bottomrule
\end{tabular}
\end{center}

\textbf{Key insight}: The dominant tropospheric systematic ($\sim$3--17\,ps)
is comparable to the 59\,ps signal in absolute terms, but it has
\emph{different} temporal structure (weather-dependent, not sidereal)
and \emph{cancels} in the multi-GNSS ratio test. With 2 years of data
and $\sim 10^5$ satellite passes, the matched filter achieves $\SNR > 100$
because the DFD template is coherent while tropospheric noise is
incoherent.

%=============================================================================
\section{Personnel, Timeline, and Budget}
%=============================================================================

\textbf{Personnel}: 1 geodesist or GNSS timing specialist (postdoc
or senior graduate student) with experience in PPP processing.

\textbf{Timeline}:
\begin{itemize}[nosep]
\item Month 1: Data download and validation. Set up PPP processing
  pipeline (GIPSY-OASIS or Bernese).
\item Months 2--3: Process GPS data, obtain clock residuals, apply
  matched filter for 59\,ps signal.
\item Month 4: Process Galileo data independently. Compute
  Galileo/GPS ratio.
\item Month 5: Systematic studies (vary tropospheric model, ionospheric
  correction, station subset). Assess robustness.
\item Month 6: Paper preparation and submission.
\end{itemize}

\textbf{Budget}:
\begin{itemize}[nosep]
\item 1 postdoc/geodesist (6 months): \$40--60K
\item Computing (standard workstation, existing PPP licenses): \$5--10K
\item Travel (1 geodesy conference): \$3--5K
\item \textbf{Total}: \$48--75K
\end{itemize}

\textbf{Funding targets}: NASA Geodesy and Geophysics program,
NSF Geophysics, ESA GNSS Science Support Centre.

%=============================================================================
\section{Success Criteria}
%=============================================================================

\begin{itemize}
\item \textbf{Detection}: 59\,ps annual signal at $\SNR > 5$
  \emph{and} Galileo/GPS ratio $R = 1.070 \pm 0.01$. This would
  be decisive evidence for DFD.
\item \textbf{Exclusion}: No 59\,ps signal at $\SNR < 2$ with
  the expected sensitivity. This constrains $\lbone < 10^{-10}$
  (improving SQMS bound by $\sim 10\times$).
\item \textbf{Systematic-limited}: Signal detected but ratio
  inconsistent with 1.070. This indicates a systematic of unknown
  origin requiring further investigation.
\end{itemize}


\newpage
%=============================================================================
%=============================================================================
\part{Proposal 3: Si$_3$N$_4$ MEMS Prototype}
\label{part:mems}
%=============================================================================
%=============================================================================

\begin{tcolorbox}[colback=green!3!white, colframe=actioncolor,
  title=\textbf{PROPOSAL 3 --- Si$_3$N$_4$ MEMS PROTOTYPE}]
\textbf{Budget}: \$200--500K\\
\textbf{Timeline}: 12--18 months\\
\textbf{Goal}: Demonstrate $Q_{\rm mech} \geq 10^9$ at 20\,mK
with metallization\\
\textbf{Key deliverable}: Characterized prototype ready for
DFD Phase~II detector integration\\
\textbf{Personnel}: 1 PI (MEMS/optomechanics) + 1 postdoc + 1
graduate student
\end{tcolorbox}

%=============================================================================
\section{Scientific Objective}
%=============================================================================

Fabricate and characterize a soft-clamped phononic crystal Si$_3$N$_4$
membrane resonator that achieves $Q_{\rm mech} \geq 10^9$ at 20\,mK
while carrying a thin metallic layer for electromagnetic coupling.
This device is the core sensor element for the DFD Phase~II MEMS
hybrid detector (W3i8), which targets $\lbone \sim 4.7\ee{-11}$ at
$1/80^{\rm th}$ the cryogenic cost of the 256-cavity Phase~III array.

The $\Qpsi$ threshold for MEMS+P11 detection is $\Qpsi \geq 4\ee{10}$
(W3i6). The mechanical Q of the MEMS sensor must exceed this to avoid
being the limiting noise source.

%=============================================================================
\section{Device Specification}
%=============================================================================

\begin{center}
\renewcommand{\arraystretch}{1.25}
\begin{tabular}{@{}p{5cm}p{4cm}p{4cm}@{}}
\toprule
\textbf{Parameter} & \textbf{Specification} & \textbf{Basis} \\
\midrule
Material & Stoichiometric Si$_3$N$_4$ (LPCVD) &
  Highest intrinsic stress \\
Film thickness & 20--50\,nm & W3i7 fab spec: 20\,nm \\
Intrinsic stress & $\sigma \geq 1.0$\,GPa &
  Dissipation dilution; Norcada spec: 1.0\,GPa \\
Membrane window & $500\,\mu\text{m} \times 500\,\mu\text{m}$ defect
  in phononic crystal & W3i7: 500\,$\mu$m defect \\
Phononic crystal & Soft-clamped design, $\geq 5$ unit cells &
  Tsaturyan \textit{et al.} (2017) design \\
Overall chip size & $5\,\text{mm} \times 5\,\text{mm}$ &
  Standard Si frame \\
Defect mode frequency & $f_0 \approx 1$\,MHz &
  W3i7 spec \\
Effective mass & $m_{\rm eff} \approx 1.2$\,ng &
  W3i7 spec; enables ng-scale psi detection \\
Target $Q$ at 300\,K & $\geq 10^8$ & Literature benchmark \\
Target $Q$ at 20\,mK & $\geq 10^9$ & DFD requirement \\
Metallization & 5--10\,nm Al or NbTi on one face &
  EM coupling; must not degrade $Q$ by $> 10\times$ \\
\bottomrule
\end{tabular}
\end{center}

%=============================================================================
\section{Fabrication Plan}
%=============================================================================

\subsection{Foundry options}

\begin{center}
\renewcommand{\arraystretch}{1.25}
\begin{tabular}{@{}p{3.5cm}p{2.5cm}p{3cm}p{4cm}@{}}
\toprule
\textbf{Vendor} & \textbf{Location} & \textbf{Capability} & \textbf{Notes} \\
\midrule
Norcada Inc. & Edmonton, Canada &
  Si$_3$N$_4$ membranes, custom MEMS, high-stress films &
  Commercial supplier of high-Q Si$_3$N$_4$ for quantum
  optomechanics. Stress-tailored films at 1.0\,GPa.
  Custom patterning available. \\
\midrule
NIST Boulder CNST & Boulder, CO, USA &
  Full nanofab including e-beam lithography &
  US government facility; academic access via user program.
  Phononic crystal patterning demonstrated. \\
\midrule
DTU Nanolab & Lyngby, Denmark &
  LPCVD Si$_3$N$_4$, e-beam lithography &
  European access; Schliesser group collaboration. \\
\midrule
Cornell CNF & Ithaca, NY, USA &
  Full MEMS fabrication &
  Academic user facility; competitive rates. \\
\bottomrule
\end{tabular}
\end{center}

\textbf{Recommended approach}: Commission Norcada for the base
Si$_3$N$_4$ membrane wafers (they have established high-stress LPCVD
processes at 1.0\,GPa). Perform phononic crystal patterning and
metallization at a university cleanroom (NIST CNST or Cornell CNF)
to maintain design flexibility during optimization.

\subsection{Fabrication process}

\begin{enumerate}
\item \textbf{Wafer preparation}: Start with 100\,mm Si(100) wafer,
  double-side polished, 200--350\,$\mu$m thick.
\item \textbf{LPCVD Si$_3$N$_4$ deposition}: Deposit 20--50\,nm
  stoichiometric Si$_3$N$_4$ on both sides. Target stress:
  $1.0 \pm 0.1$\,GPa (measured by wafer bow).
\item \textbf{Phononic crystal patterning (front side)}:
  \begin{itemize}[nosep]
  \item Spin e-beam resist (ZEP-520A or PMMA)
  \item E-beam lithography: write phononic crystal hole pattern
    (hexagonal array, pitch $\sim 200$\,$\mu$m, hole diameter
    $\sim 150$\,$\mu$m, $\geq 5$ unit cells surrounding a
    500\,$\mu$m defect region)
  \item Reactive ion etch (RIE) with CHF$_3$/O$_2$ chemistry to
    transfer pattern into Si$_3$N$_4$
  \item Strip resist
  \end{itemize}
\item \textbf{Backside release}:
  \begin{itemize}[nosep]
  \item Pattern backside Si$_3$N$_4$ with photolithography (5\,mm
    $\times$ 5\,mm window)
  \item KOH wet etch (30\% KOH, 80$^\circ$C) to remove Si substrate
    through to the front-side Si$_3$N$_4$ membrane
  \end{itemize}
\item \textbf{Metallization}: E-beam evaporate 5--10\,nm Al (or
  sputter NbTi for superconducting option) on one face through a
  shadow mask to avoid phononic crystal regions.
\item \textbf{Dicing and inspection}: Dice into individual chips;
  optical/SEM inspection for defects.
\end{enumerate}

\textbf{Yield estimate}: 10--20 chips per wafer; expect 50--70\%
yield for the phononic crystal step based on literature.

%=============================================================================
\section{Characterization Plan}
%=============================================================================

\subsection{Phase A: Room temperature (Months 1--6)}

\begin{enumerate}[nosep]
\item \textbf{Optical interferometric readout}: Fiber-based
  Fabry--Perot interferometer with membrane as one mirror.
  Measure $f_0$ and $Q$ at 300\,K in vacuum ($< 10^{-6}$\,mbar).
\item \textbf{$Q$ vs.\ membrane parameters}: Vary thickness (20, 30,
  50\,nm), phononic crystal geometry, and metallization thickness.
  Target: identify configuration with $Q(300\,\text{K}) \geq 10^8$.
\item \textbf{Stress verification}: Raman spectroscopy and/or
  bulge test to confirm $\sigma \geq 1.0$\,GPa after processing.
\end{enumerate}

\subsection{Phase B: Cryogenic characterization (Months 6--15)}

\begin{enumerate}[nosep]
\item \textbf{$Q$ vs.\ temperature}: Mount device in dilution
  refrigerator. Measure $Q$ at 4\,K, 1\,K, 100\,mK, 20\,mK.
  DFD requires $Q(20\,\text{mK}) \geq 10^9$.
\item \textbf{$Q$ vs.\ magnetic field}: Sweep $B = 0$--$1$\,T at
  20\,mK. Important for understanding metallization losses and
  compatibility with SQUID readout in Phase~II detector.
\item \textbf{Thermal noise floor}: Measure displacement noise spectral
  density. At 20\,mK with $Q = 10^9$ and $f_0 = 1$\,MHz:
  \begin{equation}
  S_x^{1/2} = \sqrt{\frac{4 k_B T}{m_{\rm eff}\,(2\pi f_0)^2\,Q\,(2\pi f_0)}}
  \sim 10^{-18}\,\text{m}/\sqrt{\text{Hz}}.
  \end{equation}
\item \textbf{Readout}: Fiber interferometric at mK temperatures.
  Use single-mode fiber passed through dilution refrigerator.
  Fiber Fabry--Perot cavity with finesse $\mathcal{F} \sim 10^4$.
\end{enumerate}

\subsection{Phase C: Metallization optimization (Months 12--18)}

\begin{enumerate}[nosep]
\item If Al metallization degrades $Q$ by $> 10\times$: try NbTi
  (superconducting, lower loss at mK).
\item If NbTi still unacceptable: try capacitive coupling geometry
  (electrode on separate chip, $\sim 1$\,$\mu$m gap) to avoid
  metallization entirely.
\item Document final $Q(T, B)$ characterization for DFD Phase~II
  integration specification.
\end{enumerate}

%=============================================================================
\section{Personnel, Timeline, and Budget}
%=============================================================================

\textbf{Personnel}:
\begin{itemize}[nosep]
\item 1 PI (MEMS/optomechanics expert, 30\% FTE): experimental design,
  student supervision, publication
\item 1 postdoc (100\% FTE): fabrication, room-temperature characterization
\item 1 graduate student (50\% FTE): cryogenic measurements
\end{itemize}

\textbf{Timeline}:
\begin{itemize}[nosep]
\item Months 1--2: Design finalization, Norcada wafer order, cleanroom
  process development
\item Months 3--6: Fabrication (2--3 wafer runs), room-temperature $Q$
  measurements
\item Months 7--12: Cryogenic cool-downs, $Q(T)$ and $Q(B)$ mapping
\item Months 13--18: Metallization optimization, final characterization,
  paper and Phase~II design document
\end{itemize}

\textbf{Budget}:
\begin{itemize}[nosep]
\item Postdoc salary (18 months): \$120--160K
\item Graduate student stipend (18 months, 50\%): \$25--35K
\item Norcada wafer orders (3 runs, custom): \$15--30K
\item Cleanroom time (NIST CNST or university fab): \$20--40K
\item Dilution refrigerator time (if shared facility): \$20--50K
\item Fiber interferometer components: \$10--20K
\item Supplies, SEM/characterization: \$10--20K
\item PI summer salary (30\% FTE, 2 months): \$15--30K
\item Travel: \$5--10K
\item \textbf{Total}: \$240--395K
\end{itemize}

\textbf{Funding targets}: DOE Office of Science (Quantum Information
Science), NSF Quantum Sensing (QuSeC), DARPA (if defense applications
emphasized).

%=============================================================================
\section{Success Criteria}
%=============================================================================

\begin{itemize}
\item \textbf{Full success}: $Q_{\rm mech} \geq 10^9$ at 20\,mK
  with metallization intact. Device ready for DFD Phase~II integration.
  Publish fabrication recipe and characterization data.
\item \textbf{Partial success}: $Q_{\rm mech} \geq 10^9$ at 20\,mK
  without metallization; metallization degrades to $Q \sim 10^8$.
  Proceed with capacitive coupling geometry for Phase~II.
\item \textbf{Minimum viable}: $Q_{\rm mech} \geq 10^8$ at 20\,mK.
  Useful for non-DFD optomechanics applications; DFD Phase~II delayed
  pending further materials R\&D.
\end{itemize}


\newpage
%=============================================================================
%=============================================================================
\part{Proposal 4: Galaxy Rotation Curve Reanalysis}
\label{part:sparc}
%=============================================================================
%=============================================================================

\begin{tcolorbox}[colback=green!3!white, colframe=actioncolor,
  title=\textbf{PROPOSAL 4 --- GALAXY ROTATION CURVE REANALYSIS}]
\textbf{Budget}: \$0 (uses only public data and standard tools)\\
\textbf{Timeline}: 3--6 months\\
\textbf{Data}: SPARC database (175 galaxies, publicly available)
and/or BIG-SPARC ($\sim$4000 galaxies, expected 2025--2026)\\
\textbf{Expected result}: $\mu_{\rm simple}$ preferred at $>5\sigma$
if DFD correct\\
\textbf{Personnel}: 1 astronomer, graduate student, or postdoc
\end{tcolorbox}

%=============================================================================
\section{Scientific Objective}
%=============================================================================

DFD uniquely predicts the ``simple'' interpolating function
$\mu(x) = x/(1+x)$, derived from the $S^3$ microsector topology
(section02\_formalism.tex, Appendix~\ref{app:mu-derivation},
Theorem~\ref{thm:mu-theorem} in the main DFD text). This is
\emph{not} a phenomenological fit --- it is a zero-parameter
prediction from the theory's mathematical structure.

The competing ``standard'' MOND interpolating function
$\mu(x) = x/\sqrt{1+x^2}$ is purely phenomenological and has no
theoretical derivation.

These two functions make distinguishable predictions for rotation
curves in the transition regime ($x \sim 1$, i.e.,
$a \sim a_0 \approx 1.2\ee{-10}$\,m/s$^2$). A Bayesian model
comparison on the SPARC rotation curve database can decisively
discriminate between them.

The key difference: at $x = 1$ (the crossover point),
$\mu_{\rm simple}(1) = 0.5$ while $\mu_{\rm standard}(1) = 1/\sqrt{2}
\approx 0.707$. This 41\% difference in the interpolation region
produces measurable differences in rotation curves of galaxies with
accelerations near $a_0$.

%=============================================================================
\section{Data Specification}
%=============================================================================

\subsection{SPARC database}

The Spitzer Photometry and Accurate Rotation Curves (SPARC) database
is available at \url{http://astroweb.cwru.edu/SPARC/}
(Lelli, McGaugh, \& Schombert, AJ 152, 157, 2016).

\textbf{Contents}:
\begin{itemize}[nosep]
\item 175 late-type galaxies with Spitzer [3.6] photometry and
  high-quality HI/H$\alpha$ rotation curves.
\item For each galaxy: rotation velocity $V_{\rm obs}(R)$, baryonic
  components $V_{\rm gas}(R)$, $V_{\rm disk}(R)$, $V_{\rm bulge}(R)$
  with uncertainties.
\item Mass-to-light ratio $\Upsilon_\star$ as a free parameter
  (or fixed at population synthesis value).
\end{itemize}

\textbf{Download}: \texttt{Rotmod\_LTG.zip} from the SPARC website
contains all rotation curve data and Newtonian mass models.

\subsection{BIG-SPARC (recommended upgrade)}

BIG-SPARC (Haubner, Lelli \textit{et al.}, arXiv:2411.13329, 2024)
expands SPARC to $\sim$4000 galaxies with homogeneous HI rotation
curves from APERTIF, ASKAP, ATCA, GMRT, MeerKAT, VLA, and WSRT,
plus WISE near-infrared photometry. When publicly released, this
should be the primary dataset as it provides $>20\times$ the
statistical power.

%=============================================================================
\section{Analysis Method}
%=============================================================================

\subsection{Step 1: Forward model}

For each galaxy $i$ with observed rotation curve
$V_{{\rm obs},i}(R_j) \pm \sigma_{ij}$, compute the predicted
rotation curve under three models:

\textbf{Model A (DFD/Simple $\mu$)}:
\begin{equation}
\mu_A(x) = \frac{x}{1+x}, \qquad x = \frac{g_{\rm bar}}{a_0},
\end{equation}
where $g_{\rm bar} = V_{\rm bar}^2/R$ is the Newtonian baryonic
acceleration. The total acceleration is:
\begin{equation}
g_{\rm obs} = \frac{g_{\rm bar}}{\mu_A(g_{\rm bar}/a_0)},
\end{equation}
so $V_{\rm pred}(R) = \sqrt{g_{\rm obs}\,R}$.

\textbf{Model B (Standard $\mu$)}:
\begin{equation}
\mu_B(x) = \frac{x}{\sqrt{1+x^2}}.
\end{equation}

\textbf{Model C (Free-form $\mu$)}:
\begin{equation}
\mu_C(x) = \frac{x}{(1 + \lambda\,x^\alpha)^{1/\alpha}},
\end{equation}
with $\alpha$ and $\lambda$ as free parameters ($\alpha = 1,
\lambda = 1$ recovers Model~A; $\alpha = 2, \lambda = 1$ recovers
Model~B).

\subsection{Step 2: Bayesian fitting}

For each galaxy and each model, the likelihood is:
\begin{equation}
\ln \mathcal{L}_i = -\frac{1}{2}\sum_j
\frac{(V_{{\rm obs},ij} - V_{{\rm pred},ij})^2}{\sigma_{ij}^2
+ \sigma_{\rm sys}^2},
\end{equation}
where $\sigma_{\rm sys}$ is an intrinsic scatter parameter
(marginalized over).

Free parameters per galaxy: $\Upsilon_\star^{\rm disk}$,
$\Upsilon_\star^{\rm bulge}$ (if applicable), distance $D$
(with Gaussian prior from measured value), inclination $i$
(with Gaussian prior).

Global parameters: $a_0$ (with tight prior around
$1.2\ee{-10}$\,m/s$^2$), and for Model~C: $\alpha$, $\lambda$.

\subsection{Step 3: Bayesian model comparison}

Compute the Bayesian evidence $\mathcal{Z}_X = \int \mathcal{L}
\cdot \pi(\theta)\,d\theta$ for each model $X \in \{A, B, C\}$
using nested sampling (\texttt{MultiNest} or \texttt{dynesty}).

The Bayes factor:
\begin{equation}
\mathcal{B}_{AB} = \frac{\mathcal{Z}_A}{\mathcal{Z}_B}
\end{equation}
provides the model comparison. $\ln \mathcal{B}_{AB} > 5$ corresponds
to decisive evidence for Model~A (Jeffreys' scale).

\subsection{Step 4: Posterior on $(\alpha, \lambda)$}

From Model~C, the joint posterior $P(\alpha, \lambda \mid \text{data})$
directly tests whether the data prefer the DFD prediction
$(\alpha = 1, \lambda = 1)$ or the standard MOND prediction
$(\alpha = 2, \lambda = 1)$. Report the 2D posterior with credible
contours.

\subsection{Expected statistical power}

With 175 SPARC galaxies, each contributing $\sim$10--30 data points
in the transition regime, the total constraining power is
$\sim 2000$--$5000$ independent measurements near $x \sim 1$.
The 41\% difference between $\mu_A(1) = 0.5$ and $\mu_B(1) = 0.707$
translates to a $\sim$20\% difference in predicted velocity at the
crossover radius. With typical velocity uncertainties of 5--10\%,
$>5\sigma$ discrimination is expected from SPARC alone.

With BIG-SPARC ($\sim$4000 galaxies), the statistical power increases
by $\sqrt{4000/175} \approx 4.8\times$, enabling $>20\sigma$
discrimination and detailed studies of environmental dependence.

%=============================================================================
\section{Systematics}
%=============================================================================

\begin{center}
\renewcommand{\arraystretch}{1.25}
\begin{tabular}{@{}p{4cm}p{3cm}p{6cm}@{}}
\toprule
\textbf{Systematic} & \textbf{Impact} & \textbf{Mitigation} \\
\midrule
Mass-to-light ratio $\Upsilon_\star$ & Degeneracy with $\mu$ shape &
  Marginalize per-galaxy; use population synthesis priors from
  Schombert \& McGaugh (2014). With Spitzer [3.6], $\Upsilon_\star$
  uncertainty is $\sim$20\% (much tighter than optical). \\
\midrule
Distance uncertainty & Affects acceleration scale &
  Include distance as a nuisance parameter with Gaussian prior
  from Tully-Fisher or tip-of-red-giant-branch measurements. \\
\midrule
Inclination uncertainty & Affects $V_{\rm obs}$ via $\sin i$ &
  Include as nuisance parameter with Gaussian prior from
  photometric axis ratio. \\
\midrule
Non-circular motions & Biases $V_{\rm obs}$ & Exclude galaxies
  with strong bars or warps (quality flag in SPARC). Use only
  ``quality 1'' and ``quality 2'' galaxies ($\sim$130 of 175). \\
\midrule
External field effect (EFE) & Modifies $\mu$ in dense environments &
  Include EFE correction using host group/cluster tidal field.
  Or restrict to isolated galaxies ($>$1\,Mpc from nearest
  massive neighbor). \\
\midrule
Gas distribution assumptions & Affects $V_{\rm gas}$ at large $R$ &
  SPARC uses tilted-ring HI models; uncertainty propagated into
  $V_{\rm bar}$ error bars. \\
\bottomrule
\end{tabular}
\end{center}

%=============================================================================
\section{Personnel, Timeline, and Budget}
%=============================================================================

\textbf{Personnel}: 1 astronomer (postdoc, graduate student, or faculty
member) with experience in galaxy dynamics and Bayesian inference.

\textbf{Timeline}:
\begin{itemize}[nosep]
\item Month 1: Download SPARC data. Implement forward model for
  Models A, B, C. Validate against published MOND fits
  (Li \textit{et al.} 2018).
\item Months 2--3: Run Bayesian fits for all 175 galaxies under each
  model. Compute Bayes factors. Produce $(\alpha, \lambda)$ posterior.
\item Month 4: Systematic studies (vary priors, galaxy quality cuts,
  $\Upsilon_\star$ assumptions, EFE corrections).
\item Months 5--6: Paper preparation. If BIG-SPARC available, extend
  analysis.
\end{itemize}

\textbf{Budget}: \$0. All data are public. Software is open-source
(\texttt{Python}, \texttt{emcee}/\texttt{dynesty}/\texttt{MultiNest}).
Computing requirements are modest (standard workstation, $\sim$days of
CPU time). The project is ideal for an existing postdoc or graduate
student seeking a high-impact publication.

\textbf{Software stack}:
\begin{itemize}[nosep]
\item \texttt{Python 3.10+}
\item \texttt{numpy}, \texttt{scipy}, \texttt{astropy}
\item \texttt{dynesty} or \texttt{pymultinest} (nested sampling)
\item \texttt{matplotlib} for visualization
\item \texttt{pygrc} (\url{https://github.com/amanmdesai/pygrc})
  for SPARC data I/O
\end{itemize}

%=============================================================================
\section{Success Criteria}
%=============================================================================

\begin{itemize}
\item \textbf{DFD confirmed}: $\ln \mathcal{B}_{AB} > 5$ (decisive
  evidence for simple $\mu$) and $(\alpha, \lambda)$ posterior centered
  on $(1, 1)$ at $> 5\sigma$. Publish as a strong test of DFD.
\item \textbf{DFD disfavored}: $\ln \mathcal{B}_{AB} < -5$ (decisive
  evidence against simple $\mu$). This would falsify the $S^3$
  microsector derivation while leaving the rest of DFD intact
  (the interpolating function could be replaced with the data-preferred
  form).
\item \textbf{Inconclusive}: $|\ln \mathcal{B}_{AB}| < 5$. Neither
  function decisively preferred. This is unlikely with full SPARC
  (expected $>5\sigma$) but possible if systematics dominate. In this
  case, BIG-SPARC data will be essential.
\end{itemize}


\newpage
%=============================================================================
%=============================================================================
\part{Cross-Cutting Considerations}
\label{part:crosscut}
%=============================================================================
%=============================================================================

%=============================================================================
\section{Proposal Interdependencies}
%=============================================================================

\begin{center}
\renewcommand{\arraystretch}{1.3}
\begin{tabular}{@{}ccp{8cm}@{}}
\toprule
\textbf{From} & \textbf{To} & \textbf{Dependency} \\
\midrule
Proposal 4 & All & If $\mu_{\rm simple}$ is disfavored at $>5\sigma$,
the $S^3$ microsector is falsified. The remaining proposals still
proceed (they test $\lambda$, not $\mu$), but the theoretical
motivation changes. \\
Proposal 3 & SQMS Phase~I & The MEMS prototype feeds directly into
the DFD Phase~II detector. SQMS Phase~I results determine whether
Phase~II proceeds with MEMS. \\
Proposal 1 & Proposal 2 & A LIGO scalar detection would dramatically
increase the priority of the GPS search (and vice versa). A joint
detection in both would be decisive. \\
Proposal 2 & Multi-GNSS & The GPS search naturally extends to the
multi-GNSS ratio test, which is the single most decisive
DFD discriminant. \\
\bottomrule
\end{tabular}
\end{center}

%=============================================================================
\section{Recommended Execution Order}
%=============================================================================

\begin{enumerate}
\item \textbf{Start immediately}: Proposal 4 (Galaxy rotation curves).
  Zero cost, 3--6 months. Any graduate student or postdoc with Python
  and astrophysics experience can begin today.

\item \textbf{Start within 1 month}: Proposal 2 (GPS 59\,ps search).
  Low cost, uses existing data. Requires a geodesist but no new
  equipment.

\item \textbf{Start within 3 months}: Proposal 1 (LIGO archival).
  Requires postdoc hiring and pipeline development. Should begin
  in parallel with Proposal 2.

\item \textbf{Start within 6 months}: Proposal 3 (Si$_3$N$_4$ MEMS).
  Longest lead time due to fabrication. Wafer orders should be placed
  as soon as funding is secured.
\end{enumerate}

%=============================================================================
\section{Relationship to SQMS Phase~I}
%=============================================================================

These four proposals are designed to \emph{complement}, not replace,
the SQMS Phase~I experiment (\$2--4M, 2027--28), which remains the
single most important DFD experiment (W3i8). The proposals here serve
three strategic functions:

\begin{enumerate}
\item \textbf{Parallel low-cost tracks}: Proposals 1, 2, and 4 can
  produce results before SQMS Phase~I begins, building the scientific
  case for DFD and attracting collaborators.

\item \textbf{Technology readiness}: Proposal 3 develops the MEMS
  sensor needed for Phase~II, which follows SQMS Phase~I.

\item \textbf{Independent confirmation}: If SQMS Phase~I detects
  a psi-signal, Proposals 1 and 2 provide independent verification
  channels using completely different physics and systematics.
\end{enumerate}

%=============================================================================
\section{Risk Assessment}
%=============================================================================

\begin{center}
\renewcommand{\arraystretch}{1.25}
\begin{tabular}{@{}clccp{5cm}@{}}
\toprule
\textbf{\#} & \textbf{Proposal} & \textbf{Risk} & \textbf{Probability} &
  \textbf{Mitigation} \\
\midrule
1 & LIGO archival & Systematic floor limits reach &
  Medium & Focus on magnetar correlation (Channel A), which has
  distinct temporal signature \\
2 & GPS 59\,ps & Tropospheric residuals mask signal &
  Low & Multi-GNSS ratio cancels atmosphere; 2\,yr integration
  averages weather \\
3 & Si$_3$N$_4$ MEMS & Metallization kills $Q$ &
  Medium & Fallback to capacitive coupling; NbTi alternative;
  useful even without metallization for other quantum sensing \\
4 & SPARC analysis & Standard $\mu$ preferred &
  Low--Medium & Would falsify $S^3$ derivation but not the DFD
  framework (replace $\mu$ with data-preferred form) \\
\bottomrule
\end{tabular}
\end{center}


%=============================================================================
\section*{Summary of Key References and Data Access Points}
%=============================================================================

\begin{itemize}[nosep]
\item \textbf{LIGO data}: GWOSC, \url{https://gwosc.org/data/}.
  O3 strain: \url{https://gwosc.org/O3/o3_details/}.
  Auxiliary channels: \url{https://gwosc.org/O3/auxiliary/}.
\item \textbf{IGS products}: \url{https://igs.org/products-access/}.
  CDDIS archive: \url{https://cddis.nasa.gov/}.
  MGEX multi-GNSS: \url{https://igs.org/mgex/products/}.
\item \textbf{Si$_3$N$_4$ MEMS vendors}: Norcada Inc.
  (\url{https://www.norcada.com/quantum-standard-devices/}),
  NIST CNST, DTU Nanolab, Cornell CNF.
\item \textbf{SPARC database}: \url{http://astroweb.cwru.edu/SPARC/}.
  BIG-SPARC: Haubner \textit{et al.}, arXiv:2411.13329 (2024).
  Zenodo mirror: \url{https://zenodo.org/records/16284118}.
\end{itemize}


\end{document}
